\paragraph{A separate book retains exact credit amounts.}
The optional \texttt{CreditLinkedBook/1} now supplies that finite projection.
It begins with a validated rational opening before acquisition, binds the
complete exact credit result, and requires explicit account roles and an
illustrative one-to-one mapping from native CU to host XCU. Cash begins at
the supplied opening amount; gross carrying amount, allowance and the three
income roles begin at zero. The original rational book remains unchanged.
Unrelated rational balances become constant expressions under the same root,
and unrelated memory remains byte-for-byte equal in its canonical encoding.

Allowance occupies a contra-asset account whose signed balance is $-L$.
For example, after the first collection in the adverse-without-write-off case,
gross carrying amount is 66 and allowance is 22, so the two signed asset
balances contribute $66-22=44$. Classifying that credit balance as a liability
would change the presentation. A write-off instead reduces gross and allowance
together; its journal has zero cash movement. Interest, impairment and catch-up
use signed income accounts, so the same roles accommodate both gains and losses.

The projection accumulates the original journals in their original order,
including zero entries. Each native observation names the exact journal prefix
that produces its balances. On 1 January 2026 in the partial-write-off case,
accrual, collection, another accrual and the improvement observation correspond
to prefixes 5, 6, 7 and 8. Sorting observations by date alone would lose that
distinction. The complete native request and result retain expected cash,
written-off portions and dated origins alongside the proposed monetary state.

Reports need only addition and the declared account classes. Let $T_A,T_L,T_E,
T_I,T_X$ denote the signed totals of asset, liability, equity, income and expense
accounts, respectively, in host currency units. The zero journal sum preserves
$T_A+T_L+T_E+T_I+T_X=0$. Presentation therefore uses
\[
 A=T_A,\qquad L=-T_L,\qquad
 P=-T_I-T_X,\qquad E=-T_E+P,\qquad A=L+E.
\]
Here $P$ is cumulative profit since the opening and $E$ includes that profit.
These relations require no ordering test on an algebraic amount. In the
constructed complement case, equipment of 17 and matching opening capital
raise assets and equity by 17 throughout, while native cash and profit are
unchanged. The period remains a cumulative interval containing the original
dates; it does not establish annual completeness or a closing policy.

All nine schedules pass on two account charts with and without that complement:
36 complete exchanges preserve 252 journals and 228 observations. The checks
certify 3,186 exact identities. Of the observations, 124 project to rational
balances accepted by the unchanged host validator, including every observation
in the three fully rational cases. This is a balance and type comparison.
The separate cash record preserves price, integral cost, fee and collection
components, without assigning operating, investing or financing classes.

That distinction has a checked consequence. Asking the existing statement
engine to explain the acquisition with no classified cash movements fails its
cash reconciliation. An absent classification cannot be replaced by a zero
movement. The optional reference consequently reports cumulative balances,
profit and unclassified cash origins; a full cash-flow statement still needs
an explicit classification policy. Full exchanges and refusals are retained
under {\small\path{data/accounting-master/}} in
{\small\path{credit-linked-book-20261025/}}. They establish a single-root
proposed sequence with no durable posting, general multi-rate composition or
accounting admission.
